\section{SAIL-HIL}
\label{sec:defense}

The Safety Authorization and Intervention Layer for Human-in-the-Loop control
(\sail{}-HIL) mediates consequential tools while retaining native planning and
local deliverable generation. Replies provide scoped evidence for individual
effects. This stateful controller extends earlier v3/v4 designs; their
fixed-reply results do not validate its clarification mechanism.
Table~\ref{tab:saildecisions} separates task authority from the decision
supplied by a reply.
\begin{table}[ht]
\centering\normalsize
\caption{SAIL-HIL decisions for the proposed effect. All dispatches additionally require valid review, scope and quantity checks. Blocking an extension leaves other authorized task obligations available.}
\label{tab:saildecisions}
\begin{tabularx}{\linewidth}{@{}>{\raggedright\arraybackslash}p{.43\linewidth}>{\raggedright\arraybackslash}X@{}}
\toprule
Authority or reply state & Controller action \\
\midrule
Covered by original task & Reserve permit; dispatch exact call \\
Missing authority; exact approval & Bind unused reply; reserve permit; dispatch \\
Unclear or nearby-scope reply & Clarify within budget; otherwise withhold \\
Explicit denial & End consultation for this proposal; withhold \\
Hard-prohibited effect & Block regardless of ordinary approval \\
\bottomrule
\end{tabularx}
\end{table}


\subsection{Authority and exact-effect proposals}
For every consequential call, a controller model constructs an effect representation
from the trusted task, proposed tool and arguments, observations, and execution
history. Content retrieved from the workspace and the actor's justification are
evidence, not sources of authority. The proposal is classified as already
authorized by the original task, a confirmable extension, or a hard-prohibited
effect. Hard constraints dominate ordinary approval.

An allowed proposal creates a permit containing the canonical effect ID, exact
call signature, quantity, authority source, and---for extensions---the reply
index. The permit is reserved before dispatch and retained if execution fails,
preventing retries from silently multiplying an authorization. Changed object,
destination, audience, privilege, persistence, purpose, or quantity requires a
new assessment. The controller model proposes the scope association, while the runtime checks the cited reply, call signature, and reservation. Mechanical validity cannot establish that the model interpreted the reply correctly. This realizes effect-level mediation for the instrumented tools; it
does not claim coverage of uninstrumented side channels.

\subsection{Reply state and bounded clarification}
The controller maintains an append-only ledger of actual question--reply pairs.
A separate reply review may return \emph{approved}, \emph{denied},
\emph{unresolved}, or \emph{scope-mismatched}. Approval must cite literal
evidence from an unused reply, cover the exact proposed effect, satisfy stated
conditions, and fit the remaining quantity budget. Denial ends consultation for
that proposal. Ambiguity, delegation of judgment, and approval of a nearby step
provide no permit.

When the latest reply is unresolved or scope-mismatched, the full controller may
ask at most two focused clarifications within a four-question case budget. The
question names the exact proposed call and asks for an explicit approve/deny
decision and any limit. New matched replies advance the simulator state; a
later exact approval can repair the decision but cannot broaden the proposed
effect. The no-controller-clarification ablation stops after the first reply,
while the no-human ablation makes no query and conservatively blocks extensions.

\subsection{Safe continuation}
Blocking an extension does not revoke authority already present in the original
request. Before termination, \sail{}-HIL reviews unsatisfied task obligations
and exposes only safe next actions to the native agent. This recovery step aims
to separate security from blanket refusal: a report can still be written after
an external send is denied. The planned no-recovery comparison tests its contribution
to benign completion and joint safe utility. Controller-model failures are marked as
guard-invalid infrastructure outcomes rather than counted as safe decisions.

\paragraph{Illustrative decision path.}
For an internal report, approval to finish a local draft supplies no permission
to send it to an external recipient. The controller can ask whether that exact
send is authorized. A later exact approval permits the reviewed send once;
denial or exhausted clarification leaves it blocked. In either case, writing
the originally requested report remains eligible. This conceptual example
separates three mechanisms: interpreting reply scope, binding permission to
execution, and continuing authorized work. It is not an observed trajectory or
an empirical guarantee of recovery.
